\section{The formal boundary: ordinary iPRA and finite specialization}
\label{sec:boundary}

\subsection{The object theory used in the report}

We use the ordinary presentation already selected by the repository
unification plan.  Its terms are built from a symbol for every primitive
recursive function.  Its formulas use intuitionistic first-order logic with
equality.  The theory \(\PRAAx\) contains the defining equations for primitive
recursive symbols, the successor axioms, and induction for quantifier-free
formulas.  In Lean the central judgment is
\[
       \Prf\;\PRAAx\;\Gamma\;\varphi,
\]
with notation \(\Gamma\vdash_{\PRAAx}\varphi\).

The checked deduction layer provides, among other rules:

\begin{itemize}
\item substitution of arbitrary terms for free variables;
\item weakening of the hypothesis list;
\item simultaneous cut, called \texttt{Prf.trans};
\item universal and existential introduction and elimination;
\item simultaneous equality substitution;
\item quantifier-free induction through \texttt{PRA.qfInd};
\item the convenient open induction wrapper \texttt{LeanProofs.IPRA.ind};
\item instantiation of closed derivations at arbitrary terms through
      \texttt{LeanProofs.IPRA.inst}.
\end{itemize}

The framework never reasons by saying that a primitive recursive function has
the intended standard value and then applying a semantic completeness
principle.  Instead, every use of a program is mediated by its defining
equations and an object derivation.

\subsection{Four levels that must not be conflated}

For a fixed algorithm \(f\), formula \(A\), and certificate \(c\), there are
four different statements.

\begin{enumerate}
\item \textbf{Host computation:} a Python or Lean evaluator returns a value.
\item \textbf{Host theorem:} Lean proves a proposition about the standard
      natural-number interpretation of \(f\).
\item \textbf{Object equation:} \(\ipra\) derives the defining or correctness
      equation for the term denoting \(f\).
\item \textbf{Object conclusion:} \(\ipra\) derives \(A\), perhaps after
      eliminating the certificate and discharging residual assumptions.
\end{enumerate}

Levels 1 and 2 are useful for discovering invariants and regression testing,
but only levels 3 and 4 prove the requested theorem.  FIBER metadata records
all four separately.

\subsection{Why the compiler is specialized}

One might try to internalize the entire framework by adding a
primitive-recursive universal evaluator \(U(e,x)\) for codes of arbitrary
primitive-recursive functions.  This cannot work.

\begin{theorem}[No primitive-recursive universal evaluator]
There is no binary primitive-recursive function \(U\) such that for every
unary primitive-recursive function \(f\) there is an index \(e\) satisfying
\(U(e,x)=f(x)\) for every \(x\).
\end{theorem}

\begin{proof}
If such \(U\) existed, the diagonal function
\[
             g(x)=U(x,x)+1
\]
would be primitive recursive.  Hence \(g(x)=U(e,x)\) for some \(e\).  Setting
\(x=e\) gives \(U(e,e)+1=U(e,e)\), impossible.  The argument is constructive
and uses only closure of primitive recursive functions under composition and
successor.
\end{proof}

Consequently a framework that quantifies internally over all primitive
recursive source programs is asking for the wrong interface.  FIBER uses
\emph{metalevel specialization}: a finite source description is compiled to
specific \(\mathsf{PR}\) terms, specific formulas, and a specific proof term.
For a bounded dynamic trace, the trace checker itself is a fixed primitive
recursive symbol.  The generated proof is then checked by Lean's kernel.

\subsection{Finite-use discipline}

A proof of a fixed finite-query conclusion can inspect only finitely much
source data.  FIBER makes that finite use explicit through a demand map.  This
is not merely an implementation optimization.  It is the logical reason that
objects such as Cauchy reals, analytic functions, compact spaces, and quotient
objects can be handled in weak arithmetic: the theorem does not need the
completed infinitary object; it needs a finite observation and a proof that
this observation is stable under the chosen representation.

The discipline has two forms.

\begin{description}
\item[Static finite use.] The compiler knows a finite list of source queries
      from the target query.
\item[Adaptive finite use.] A bounded interaction program chooses its next
      query from previous answers, but a primitive-recursive fuel or a finite
      monotone closure bounds the complete trace.
\end{description}

Both forms produce one finite master certificate.  No appeal to countable
choice or bar induction is needed for the compiled theorem.

\subsection{Portable versus backend-specific statements}

The canonical portable conclusion is an ordinary \(\Prf\;\PRAAx\) sequent.
When an application is stated in the Foundation syntax, the checked bridge
transports the resulting open or closed sequent.  This separation avoids two
common mistakes:

\begin{enumerate}
\item proving the right mathematics in a nonportable host representation and
      never obtaining an object theorem;
\item rebuilding arithmetic inside every application-specific Foundation
      development.
\end{enumerate}

The framework therefore compiles first to the ordinary calculus, imports the
checked arithmetic library, and translates only at the application boundary.
